% Focus: Present the contact structure and the horizontal differential calculus as the local geometric layer.

\subsection{The arithmetic expression contact form}\label{subsec:contact-form}
% Focus: Define the contact manifold and form; prove the contact condition and connect to propagation directions.

The ACS of Section~\ref{subsec:acs-definition} records additive and multiplicative charges but forgets the current assignment value. To localize torsion while retaining the assignment, we pass to a three-dimensional space with coordinates $(u,v,a)$, where $u$ and $v$ play the local additive and multiplicative roles and $a$ records the current assignment.

We work on the manifold
\[
\mathcal{C}\coloneqq \mathbb{R}^3_{(u,v,a)}.
\]
Define the $1$-forms
\begin{equation}\label{eq:contact-forms}
\omega\coloneqq \mu\,du+\lambda a\,dv,
\qquad
\alpha\coloneqq da-\omega.
\end{equation}
For a tangent vector $X=x_u\partial_u+x_v\partial_v+x_a\partial_a$ one computes
\[
\alpha(X)=x_a-\mu x_u-\lambda a x_v.
\]
Thus $\alpha(X)=0$ if and only if $x_a=\mu x_u+\lambda a x_v$, which is exactly the infinitesimal propagation law from Proposition~\ref{prop:flow-equation} in Pfaff form.

\begin{proposition}\label{prop:contact-form}
If $\mu\lambda\neq 0$, then $\alpha$ is a contact form on $\mathcal{C}$.
\end{proposition}

\begin{proof}
Since $d\omega=\lambda\,da\wedge dv$, we have $d\alpha=-\lambda\,da\wedge dv$ and therefore
\[
\alpha\wedge d\alpha
=(da-\mu\,du-\lambda a\,dv)\wedge(-\lambda\,da\wedge dv)
=\mu\lambda\,du\wedge da\wedge dv,
\]
which is nonzero everywhere when $\mu\lambda\neq 0$.
\end{proof}

The arithmetic form \eqref{eq:contact-forms} is a coordinatization of the standard local contact model. After rescaling $\tilde u=\mu u$ and $\tilde v=\lambda v$ and setting $x=\tilde v$, $y=a$, $z=a-\tilde u$, one obtains $\alpha=dz-y\,dx$ \cite{Darboux1882Pfaff,Geiges2008IntroContact}. The key point is not the existence of a contact structure, but the distinguished horizontal directions singled out by addition and multiplication.

\subsection{Horizontal distribution and local non-integrability}\label{subsec:contact-horizontal}
% Focus: Define the horizontal distribution; compute brackets and interpret local commutator curvature.

The \emph{horizontal distribution} is the rank-$2$ distribution
\[
\mathcal{H}\coloneqq \ker\alpha\subset T\mathcal{C}.
\]
The two basic horizontal lifts of the coordinate fields on the $(u,v)$-plane are
\begin{equation}\label{eq:contact-horizontal-basis}
D_u\coloneqq \partial_u+\mu\,\partial_a,
\qquad
D_v\coloneqq \partial_v+\lambda a\,\partial_a.
\end{equation}
They satisfy $\alpha(D_u)=\alpha(D_v)=0$, hence $\mathcal{H}=\mathrm{span}\{D_u,D_v\}$. A horizontal curve $\gamma(s)$ is characterized by $\dot\gamma(s)\in\mathcal{H}$, equivalently by the Pfaff constraint $\alpha(\dot\gamma)=0$.

\begin{proposition}\label{prop:contact-bracket}
The horizontal distribution is not integrable. More precisely,
\begin{equation}\label{eq:contact-bracket}
[D_u,D_v]=\mu\lambda\,\partial_a.
\end{equation}
\end{proposition}

\begin{proof}
Compute directly:
\[
[D_u,D_v]=[\partial_u+\mu\partial_a,\partial_v+\lambda a\partial_a]=\mu\lambda\,\partial_a.
\]
\end{proof}

Equation~\eqref{eq:contact-bracket} is the local differential shadow of arithmetic torsion: an infinitesimal $u$-move followed by an infinitesimal $v$-move differs from the reversed order by a purely vertical drift. At finite scale this drift carries the same exponential weighting seen globally in the ACS.

\subsubsection*{A finite commutator square}
Let $\Phi_u^h$ and $\Phi_v^k$ denote the time-$h$ and time-$k$ flows of $D_u$ and $D_v$. They are
\[
\Phi_u^h(u,v,a)=(u+h,v,a+\mu h),
\qquad
\Phi_v^k(u,v,a)=(u,v+k,e^{\lambda k}a).
\]
Comparing the two two-step histories with the same base endpoint,
\[
p_{uv}\coloneqq \Phi_v^k\!\circ\Phi_u^h(u,v,a),
\qquad
p_{vu}\coloneqq \Phi_u^h\!\circ\Phi_v^k(u,v,a),
\]
one finds
\begin{equation}\label{eq:contact-finite-torsion}
a(p_{uv})-a(p_{vu})=\mu h\bigl(e^{\lambda k}-1\bigr),
\end{equation}
whose first-order part is $\mu\lambda hk$.

\subsubsection*{A transmission identity to the ACS weighting}
Identify the local charge plane $(A,M)$ of Section~\ref{subsec:acs-definition} with the $(u,v)$-plane by
\[
A=u,
\qquad
M=\lambda v.
\]
Then the ACS $1$-form $\eta=e^M\,dA$ pulls back to
\[
\eta_{\mathrm{loc}}\coloneqq e^{\lambda v}\,du.
\]
Its exterior derivative is the weighted area density
\begin{equation}\label{eq:contact-eta-local}
d\eta_{\mathrm{loc}}
=\lambda e^{\lambda v}\,dv\wedge du.
\end{equation}

\begin{proposition}[Local ACS weighting from contact propagation]\label{prop:contact-acs-bridge}
Let $R=[u,u+h]\times[v,v+k]$ be an oriented rectangle in the $(u,v)$-plane and let $p_{uv},p_{vu}$ be the two corresponding two-step histories based at $(u,v,a)$ as above. Then
\begin{equation}\label{eq:contact-acs-bridge}
a(p_{uv})-a(p_{vu})
=\iint_R \mu\,d\eta_{\mathrm{loc}}
=\oint_{\partial R}\mu\,\eta_{\mathrm{loc}}.
\end{equation}
\end{proposition}

\begin{proof}
Using \eqref{eq:contact-eta-local},
\[
\iint_R \mu\,d\eta_{\mathrm{loc}}
=\int_{u}^{u+h}\int_{v}^{v+k}\mu\lambda e^{\lambda v'}\,dv'\,du'
=\mu h\bigl(e^{\lambda k}-1\bigr),
\]
which is exactly \eqref{eq:contact-finite-torsion}. Stokes' theorem in the $(u,v)$-plane gives the boundary expression.
\end{proof}

\subsection{The horizontal differential \texorpdfstring{$\delta$}{delta}}\label{subsec:contact-delta}
% Focus: Define the horizontal calculus; state basic identities; clarify dependence on the chosen structure.

The horizontal distribution comes with a natural differential operator obtained by projecting the de Rham differential to the horizontal directions.

\begin{definition}\label{def:delta}
For a smooth scalar field $F(u,v,a)$ on $\mathcal{C}$, its \emph{horizontal differential} is
\begin{equation}\label{eq:delta-def}
\delta F\coloneqq (D_uF)\,du+(D_vF)\,dv.
\end{equation}
\end{definition}

\begin{lemma}\label{lemma:delta-projection}
For any smooth $F$,
\begin{equation}\label{eq:delta-projection}
\delta F=dF-(\partial_aF)\alpha.
\end{equation}
In particular, $\delta a=\omega$ and $\delta u=du$, $\delta v=dv$.
\end{lemma}

\begin{proof}
Expand $dF=(\partial_uF)\,du+(\partial_vF)\,dv+(\partial_aF)\,da$ and use $da=\alpha+\omega$ with \eqref{eq:contact-forms}. Then
\[
dF-(\partial_aF)\alpha
=(\partial_uF+\mu\partial_aF)\,du+(\partial_vF+\lambda a\,\partial_aF)\,dv
=(D_uF)\,du+(D_vF)\,dv.
\]
The special cases follow by taking $F=a,u,v$.
\end{proof}

The operator $\delta$ is not nilpotent: its failure to square to zero is precisely the commutator curvature of the horizontal distribution.

\begin{proposition}\label{prop:delta-square}
For any smooth $F$,
\begin{equation}\label{eq:delta-square}
\delta^2F=\mu\lambda(\partial_aF)\,du\wedge dv.
\end{equation}
In particular, $\delta^2a=\mu\lambda\,du\wedge dv$.
\end{proposition}

\begin{proof}
By definition, $\delta^2F=([D_u,D_v]F)\,du\wedge dv$. Using \eqref{eq:contact-bracket} gives \eqref{eq:delta-square}.
\end{proof}

This local $2$-form density $\mu\lambda\,du\wedge dv$ will later appear as the curvature term in the horizontal holomorphic and twisted harmonic operators.

\subsection{Degenerate limits and the singular model}\label{subsec:contact-degenerate}
% Focus: Analyze the cases $\mu=0$ or $\lambda=0$; describe the resulting precontact/foliated structures and what survives.

The contact condition in Proposition~\ref{prop:contact-form} fails when one of the propagation intensities vanishes. This is consistent with arithmetic intuition: if one generator is missing, the commutator defect disappears.

\subsubsection*{The additive-only limit \texorpdfstring{$\lambda=0$}{lambda=0}}
If $\lambda=0$, then $\alpha=da-\mu\,du$ and $d\alpha=0$, so $\alpha\wedge d\alpha=0$. The horizontal distribution is spanned by $\partial_v$ and $\partial_u+\mu\partial_a$ and is integrable. In this limit, $D_v=\partial_v$ and $[D_u,D_v]=0$, hence $\delta^2=0$.

\subsubsection*{The multiplicative-only limit \texorpdfstring{$\mu=0$}{mu=0}}
If $\mu=0$, then $\alpha=da-\lambda a\,dv$ and again $\alpha\wedge d\alpha=0$. The horizontal distribution is spanned by $\partial_u$ and $\partial_v+\lambda a\partial_a$ and is integrable. Again $[D_u,D_v]=0$ and $\delta^2=0$.

These degenerate limits are precontact/foliated structures rather than contact structures. They reflect the purely additive or purely multiplicative regimes where torsion vanishes. This degeneration should not be confused with the singular AES model $\mathfrak{E}_1$ of Section~\ref{subsec:aes-models}, where $\mu\lambda\neq 0$ but the assignment has a geometric singularity at its zero locus.
